\documentclass[11pt]{article}
\usepackage[T1]{fontenc}
\usepackage{lmodern,amsmath,amsthm,amssymb,microtype,xcolor,fancyhdr}
\usepackage[letterpaper,margin=1.03in,top=.9in,bottom=.92in,headheight=14pt]{geometry}
\usepackage[colorlinks=true,linkcolor=teal!50!black,citecolor=teal!50!black,urlcolor=teal!60!black]{hyperref}
\definecolor{ink}{HTML}{18394B}
\newtheorem{theorem}{Theorem}
\newtheorem{lemma}[theorem]{Lemma}
\newtheorem{corollary}[theorem]{Corollary}
\newtheorem{proposition}[theorem]{Proposition}
\DeclareMathOperator{\Log}{Log}
\newcommand{\model}{X_{\chi,h,A}}
\newcommand{\repourl}{https://github.com/lsernau/grh-lean/blob/72fead102f663d5c7a407964a13dc0cf9d977207/grh}
\pagestyle{fancy}\fancyhf{}\fancyhead[L]{\small\color{ink}Sampled Li coefficients}\fancyhead[R]{\small Research note}\fancyfoot[C]{\thepage}
\setlength{\emergencystretch}{2em}
\hypersetup{pdftitle={Identifying a Dirichlet character from filtered Li coefficients},pdfauthor={OpenAI}}
\begin{document}
\thispagestyle{plain}
{\noindent\sffamily\small\color{ink} ANALYTIC NUMBER THEORY\hfill 26 NOVEMBER 2026\par}
\vspace{.8em}
{\noindent\LARGE\bfseries\color{ink}Identifying a Dirichlet character\\[2pt] from filtered Li coefficients\par}
\vspace{.7em}
\noindent\textbf{Abstract.}
An unknown primitive Dirichlet character, a positive integer sampling stride, and a normalized finite sequence of complex weights are jointly determined by filtered samples of generalized Li coefficients. Observations may occupy widely separated intervals, with a polynomial of bounded degree fitted independently on each interval and arbitrary errors on a vanishing fraction of its indices. Distinct models differ by a fixed amount on a positive proportion of every sufficiently long interval after every such fit. We recover integer zero multiplicities by rounding spectral weights near their accumulation point. Two terms of the zero count then reveal the stride and the conductor. Dirichlet GRH gives a direct Fourier interpretation. A proof using reflection of the zeros establishes uniqueness independently of GRH.

\section{The observation and what it determines}
Let $\chi$ be a primitive Dirichlet character of conductor $f$, including $f=1$ for the Riemann zeta function. Write $Z_\chi$ for the multiset of its nontrivial zeros. Use the complex Li coefficients in the convention of Lagarias \cite[Lemma 2.3]{lagarias},
\begin{equation}
 u_\rho=1-1/\rho,\qquad
 \lambda_\chi(k)=\lim_{T\to\infty}\sum_{\substack{\rho\in Z_\chi\\|\rho|\le T}}(1-u_\rho^k),
 \qquad k=0,1,2,\ldots. \label{eq:li}
\end{equation}
The limit uses multiplicities. Let $h\ge1$ be an integer and $A(z)=\sum_{j=0}^J a_jz^j$ a complex polynomial with $A(1)=1$. The observation model is
\begin{equation}
 \model(n)=\sum_{j=0}^J a_j\lambda_\chi(hn+j),\qquad n=0,1,2,\ldots. \label{eq:model}
\end{equation}
Zero coefficients are allowed. Thus $A(z)=z^r$ encodes an unknown starting offset $r$, and several nonzero coefficients encode a weighted combination of offsets. Equality of filters means equality of polynomials.

\begin{theorem}[Uniform separation]\label{thm:main}
For two distinct triples $(\chi,h,A)$ and $(\psi,k,B)$ as above and any fixed integer $D\ge0$, there exist $\epsilon,\eta>0$ and $N_0$ such that for every interval $I$ of at least $N_0$ consecutive nonnegative integers and every complex polynomial $P$ of degree at most $D$,
\[
 \#\{n\in I:\ |X_{\chi,h,A}(n)-X_{\psi,k,B}(n)-P(n)|>\epsilon\}
 \ \ge\ \eta |I|.
\]
The constants may depend on the two triples and $D$.
\end{theorem}

\begin{corollary}[Incomplete observations]\label{cor:id}
Suppose $|I_j|\to\infty$ and data on these intervals admit a representation $\model(n)+P_j(n)+e_j(n)$, with $\sup_j\deg P_j<\infty$. For every $t>0$ suppose the fraction of $n\in I_j$ with $|e_j(n)|>t$ tends to zero. Then the triple $(\chi,h,A)$ is unique.
\end{corollary}
\begin{proof}
Two representations give a polynomial difference of bounded degree. Outside the union of the two error sets at level $\epsilon/2$, the residual difference has magnitude at most $\epsilon$. Theorem~\ref{thm:main} gives a contradiction for distinct triples.
\end{proof}
For example, take $I_j=[2^{j^2},2^{j^2}+j^2)\cap\mathbb Z$, whose union has natural density zero. Samples of $\lambda_\chi(11n+2)$ and $\lambda_\chi(11n+5)$ mixed with weights $1/4$ and $3/4$ determine the character, the integer $11$, both offsets and both weights, even with an arbitrary quadratic background on each $I_j$ and arbitrary replacements at square indices. Each sufficiently late window contains at most one square. A competing representation can have any finite filter degree.

\section{How the two supplied sources enter}
The RH article \cite{wiki} presents the critical-line formulation and Li's positivity criterion \cite{li}. The substitution in \eqref{eq:li} converts the critical line into the unit circle because $|\rho-1|=|\rho|$ there. Keeping each complex phase gives data about individual zeros. How much of that information survives the observation in \eqref{eq:model}?

At the pinned revision of \cite{repo}, \href{\repourl/src/Statements.lean}{\texttt{Statements.lean}} defines the open-strip and half-plane forms using Mathlib's Dirichlet $L$-function. In \href{\repourl/src/Main.lean}{\texttt{Main.lean}}, \texttt{LFunction\_eq\_zero\_re} places every open-strip zero on $\Re s=1/2$ and \texttt{LFunction\_ne\_zero\_of\_half\_lt\_re} gives right-half-plane nonvanishing with the principal-character exception at $s=1$. The \href{\repourl/GUIDE.md}{reading guide} traces the latter to a one-sided bound for a weighted prime-sum remainder. Short integral forms, a spanning result, and a lower bound on determinants supply the main inequality. The Lean declarations supply the circle geometry in Section~\ref{sec:fourier}. The independent argument in Sections~\ref{sec:diff}--\ref{sec:finish} only requires classical properties of Dirichlet $L$-functions.

Lagarias gives an interpolating function for the full Li sequence and, under the critical-line hypothesis, a distributional Fourier transform with frequencies accumulating only at zero \cite[Theorem 7.1]{lagarias}. His Theorem 1.1 also gives the leading $n\log n$ term under that hypothesis. Banks distinguishes distinct primitive characters of the same modulus from their zeros in short height intervals \cite{banks}. Our result concerns joint identification after unknown filtering, subsampling, polynomial fitting and sparse corruption. A targeted literature search located these predecessors and no statement of the joint theorem. Historical priority remains open. We inspected the source declarations and guide. The new result has a mathematical proof here and finite numerical checks in the archive. The published machine checks pertain to the upstream developments \cite{repo,lean}.

\section{Differencing and observable exponential terms}\label{sec:diff}
We use three standard facts. The completed entire $L$-function has order one with zero multiset $Z_\chi$. For conductor one we remove the poles of the completion by multiplying by $s(s-1)$. The functional equation and conjugation give the multiplicity-preserving reflection $\rho\mapsto 1-\overline\rho$ within $Z_\chi$. Finally the count with both signs of the ordinate is
\begin{equation}
 N_\chi(T)=\#\{\rho\in Z_\chi:|\Im\rho|\le T\}
 =\frac{T}{\pi}\log\frac{fT}{2\pi e}+O_\chi(\log(2+T)). \label{eq:rv}
\end{equation}
See \cite[Theorem 2.1]{lagarias} and \cite{bennett}. In particular, $\sum_\rho|\rho|^{-2}<\infty$ by partial summation, and a zero at height $T$ has multiplicity $O_\chi(\log(2+T))$ by differencing the count across a unit neighborhood.

Let $\Delta x(n)=x(n+1)-x(n)$ and choose $d=\max(2,D+1)$. Applying $\Delta^d$ to \eqref{eq:model} gives
\begin{align}
 \Delta^d \model(n)&=\sum_w b_{\chi,h,A}(w)\,w^n, \label{eq:exp}\\
 b_{\chi,h,A}(w)&=-(w-1)^d
   \sum_{\substack{\rho\in Z_\chi\\u_\rho^h=w}} A(u_\rho). \label{eq:atom}
\end{align}
Absent coefficients are zero. First apply the difference operator to the common truncation in \eqref{eq:li}. The constant term disappears and each power picks up $(u_\rho^h-1)^d$. Since $u_\rho^h-1=O_h(1/|\rho|)$ and the nodes $u_\rho$ are bounded, the resulting coefficients have summable absolute values. For fixed $n$ the powers $w^n$ are bounded over the nodes, so we can remove the cutoff and group terms. The only possible accumulation point of the $w$ is $1$. Contributions with $w=1$ vanish.

\begin{lemma}[Detection on translated intervals]\label{lem:shell}
Let $s(n)=\sum_w c_w w^n$ have distinct bounded nodes accumulating only at $1$ and $M=\sum_w|c_w|<\infty$. Suppose some $c_w\ne0$ for $|w|\ge1$. Choose a nonzero coefficient $c_v$ of maximal node modulus $R=|v|\ge1$. Then
\[
 \lim_{N\to\infty}\frac1N\sum_{n=L}^{L+N-1}s(n)v^{-n}=c_v
\]
uniformly for nonnegative integers $L$. On every sufficiently long interval at least the fraction $|c_v|/(4M)$ of its indices satisfy $|s(n)|\ge |c_v|R^n/2$.
\end{lemma}
\begin{proof}
A maximum exists since for every $\delta>0$ there are finitely many nodes with $|w|\ge1+\delta$. In the absence of nonzero exterior coefficients one takes a unit node. For each $w\ne v$ with $c_w\ne0$ the geometric average of $(w/v)^n$ is at most $2/(N|1-w/v|)$ in absolute value, uniformly in $L$. It is also at most one. Sum first over finitely many nodes, then bound the omitted tail by its $\ell^1$ norm. For large $N$ the modulus of the average is at least $3|c_v|/4$. Every summand $s(n)v^{-n}$ has magnitude at most $M$, so the fraction with magnitude at least $|c_v|/2$ is at least $|c_v|/(4M)$.
\end{proof}

\section{Recovering multiplicities, stride and conductor}\label{sec:decode}
\begin{proposition}\label{prop:decode}
If the coefficients in \eqref{eq:atom} for two triples agree at every $|w|\ge1$, then their strides and conductors agree. Their filters are equal and their nontrivial zero multisets differ at only finitely many points, counted with multiplicity.
\end{proposition}
\begin{proof}
\emph{One preimage near the accumulation point.}
For $w$ sufficiently close to $1$ with $w\ne1$, an occupied node of one model has the single distinct preimage
\[
 u=\exp((\Log w)/h).
\]
Here $\Log$ is principal. All but finitely many $u_\rho$ lie in an arbitrarily small neighborhood of $1$ on which $u\mapsto u^h$ is injective with this inverse. Shrinking the punctured neighborhood of $w=1$ excludes the images of the finitely many other nodes. An image exactly equal to $1$ is already omitted. The same argument applies to the competing model. Also $A$ is nonzero near $1$.

\emph{Integer recovery.}
For an occupied exterior or unit node in this neighborhood, put $m=m_\chi(\rho)$. The observable
\[
 H(w)=-\frac{b_{\chi,h,A}(w)}{(w-1)^d}=m A(u)
\]
has magnitude within $o(1)$ of $m$. Indeed
\[
 \big|\,|H(w)|-m\,\big|\le m|A(u)-1|
 =O_{\chi,A}\!\left(\frac{\log(2+|\Im\rho|)}{|\rho|}\right)\longrightarrow0 .
\]
Thus nearest-integer rounding recovers the full multiplicity, even for repeated zeros and complex weights. In particular two models with the stated coefficient equality have the same occupied exterior and unit nodes with the same rounded multiplicities near $1$.

\emph{Reflection.}
For $\rho'=1-\bar\rho$ one has $u_{\rho'}=1/\bar u_\rho$, and hence $w'=1/\bar w$. Reflection supplies the unweighted interior nodes and their multiplicities from the exterior ones. Choose a small neighborhood $|\Log w|<\delta_0$ with $\delta_0<1$, invariant under this reflection. All its occupied nodes except $1$ are now known with multiplicity. Reflection is applied after rounding, since complex filter values need not respect conjugation.

\emph{Counting at the puncture.}
Write $C(\delta)$ for the number with multiplicity of these nodes with $\delta\le |\Log w|<\delta_0$. For high zeros $\rho=\beta+i\gamma$, uniformly in $0<\beta<1$,
\[
 |\Log w|=h\left|\Log(1-1/\rho)\right|
 =\frac{h}{|\gamma|}+O_h(|\gamma|^{-2}).
\]
Consequently the height cutoff differs from $h/\delta$ by at most a bounded amount, in the sense of upper and lower inclusions. The omitted low zeros give a bounded contribution. The bounded height band at the moving cutoff contains $O_{\chi,h}(\log(1/\delta))$ zeros by \eqref{eq:rv}. Thus
\begin{equation}
 C(\delta)=\frac{h}{\pi\delta}
 \left(\log\frac1\delta+\log\frac{fh}{2\pi e}\right)
 +O_{\chi,h,\delta_0}\!\left(\log\frac1\delta\right). \label{eq:count}
\end{equation}
The two displayed coefficients recover $h$ and $f$, in that order. Changing the small fixed outer radius changes the count by $O(1)$.

\emph{Recovering the filter and the zeros.}
With $h$ common, the unfolded $u$ are identical for the two models throughout the punctured neighborhood. Equality of $H(w)$ and of the nonzero multiplicities gives $A(u)=B(u)$ on the exterior and unit part. There are infinitely many such $u$ by \eqref{eq:rv} and reflection. A polynomial vanishing at infinitely many distinct points is zero, so $A=B$. Every sufficiently high zero, on either side, is now common with its multiplicity by $\rho=1/(1-u)$.
\end{proof}

\section{From the zeros to the character}\label{sec:finish}
\begin{lemma}\label{lem:char}
If two primitive Dirichlet $L$-functions have the same nontrivial zeros with multiplicities outside a bounded set, their characters agree.
\end{lemma}
\begin{proof}
Write $F_\chi(s)=p_\chi(s)(f/\pi)^{(s+a)/2}\Gamma((s+a)/2)L(s,\chi)$, with $a\in\{0,1\}$ the parity and $p_\chi(s)=s(s-1)$ in conductor one, $1$ otherwise. Define $F_\psi$ similarly. Hadamard factorization of these order-one entire functions yields $F_\chi/F_\psi=e^{bs+c}R(s)$ for a rational function $R$. Common canonical factors cancel and the finitely many leftover exponential factors have linear exponents. Put $H(s)=L(s,\chi)/L(s,\psi)$. Forming the quotient at $s+2$ and $s$ makes the exponential factor constant and both gamma shifts rational by $\Gamma(z+1)=z\Gamma(z)$. It follows that $H(s+2)/H(s)$ is rational, including its constant prefactor. For positive real $\sigma\to\infty$ the Dirichlet series give $H(\sigma)=1+O(2^{-\sigma})$. A rational function differing from one by $O(2^{-\sigma})$ on this ray equals one identically, by its Laurent expansion at infinity. Therefore $H$ is $2$-periodic. For $\Re s>1$, $H(s)=\lim_{j\to\infty}H(s+2j)=1$. Uniqueness of absolutely convergent Dirichlet series gives $\chi(n)=\psi(n)$ for every positive integer $n$. Primitivity identifies their primitive characters and conductors.
\end{proof}

\begin{proof}[Proof of Theorem~\ref{thm:main}]
For distinct triples, Proposition~\ref{prop:decode} and Lemma~\ref{lem:char} imply that the coefficient difference $c_w$ is nonzero at some $|w|\ge1$. Apply Lemma~\ref{lem:shell} to $s=\Delta^d(X_{\chi,h,A}-X_{\psi,k,B})$. Write $v,R,M$ as in that lemma. Every polynomial $P$ of degree at most $D$ is killed by $\Delta^d$, and the absolute values of the binomial difference coefficients sum to $2^d$. If $|s(n)|\ge |c_v|R^n/2$, some $n+j$, $0\le j\le d$, satisfies
\[
 |X_{\chi,h,A}(n+j)-X_{\psi,k,B}(n+j)-P(n+j)|
 \ge |c_v|/2^{d+1}.
\]
Apply the fraction bound to the interval obtained by removing the last $d$ indices of $I$. Each witness index is used at most $d+1$ times. For sufficiently long intervals we may take $\epsilon=|c_v|/2^{d+2}$ and $\eta=|c_v|/[8M(d+1)]$, uniformly in the coefficients of $P$.
\end{proof}

\section{Reading the spectrum under Dirichlet GRH}\label{sec:fourier}
The Lean theorem gives $|w|=1$ for every zero. The sequence $\Delta^d\model$ is then bounded and has an absolutely convergent Fourier series. Every individual coefficient is recovered by
\begin{equation}
 b(w)=\lim_{N\to\infty}\frac1N\sum_{n=L}^{L+N-1}
          \Delta^d\model(n)\,\bar w^{\,n},\qquad |w|=1, \label{eq:mean}
\end{equation}
uniformly in $L$. The proof is the geometric average used in Lemma~\ref{lem:shell}, including for absent nodes. All nodes can be read from these means. Equation~\eqref{eq:count} gives explicit limits for the stride and conductor.
\begin{align*}
 h&=\lim_{\delta\downarrow0}\frac{\pi\delta C(\delta)}{\log(1/\delta)},\\
 f&=\frac{2\pi e}{h}\exp\left(\lim_{\delta\downarrow0}
     \left[\frac{\pi\delta}{h}C(\delta)-\log(1/\delta)\right]\right).
\end{align*}
Finitely many wrong counts outside the neighborhood in which rounding is exact leave the limits unchanged.

If the measured data contain unbounded outliers, a raw average in \eqref{eq:mean} can fail. Take $d\ge2$ above the background degree and let $T_K(z)$ be radial projection onto the closed disk of radius $K$. In each observed window use only indices where the whole difference stencil is available. For fixed $K$, replacing $T_K(\Delta^d\model)$ by $T_K(\Delta^d\text{data})$ changes the demodulated mean by $o(1)$ along the windows. Outside at most $d+1$ translates of the bad set the projection error is at most $2^dt$ by nonexpansiveness, and on that set it is at most $2K$. Send the window length to infinity and then $t$ to zero. For every $K>\sum_w|b(w)|$ the clean projection is the identity, so this procedure gives \eqref{eq:mean}. This determination uses infinite data. Frequencies can lie arbitrarily close together, and division by $(w-1)^d$ amplifies finite-sample errors near $1$.

\begin{corollary}[Real observations]\label{cor:real}
With real filter coefficients, observations of $\Re X_{\chi,h,A}$ under the conditions of Corollary~\ref{cor:id} determine $h$, $A$ and the unordered pair $\{\chi,\bar\chi\}$. Complex conjugation is exactly the ambiguity in the character.
\end{corollary}
\begin{proof}
Conjugating the zeros in the defining truncation gives $\lambda_{\bar\chi}(n)=\overline{\lambda_\chi(n)}$. Hence twice the real observation is the same filtered Li sum for the union, with multiplicities, of $Z_\chi$ and $Z_{\bar\chi}$. This union satisfies reflection, the same logarithmic multiplicity bound and exactly twice the count in \eqref{eq:rv}, since conjugation preserves absolute height. The argument through Proposition~\ref{prop:decode} applies with the two terms of \eqref{eq:count} doubled. In detail, Lemma~\ref{lem:shell} and the difference stencil force equality of all exterior and unit coefficients for indistinguishable observations, integer rounding recovers multiplicities for the union, and its count gives the common stride and conductor. Equality of the recovered polynomial values gives the common filter.

The two resulting products of completed functions agree in their zeros up to a finite divisor. The proof of Lemma~\ref{lem:char} applies to the two gamma factors in each product and gives $L(s,\chi)L(s,\bar\chi)=L(s,\psi)L(s,\bar\psi)$. Equality of the Dirichlet coefficient at each prime $p\nmid f$ gives $(\chi+\bar\chi)(p)=(\psi+\bar\psi)(p)$. Every reduced residue class modulo $f$ contains a prime, so the character sums are equal on the finite group of units. Orthogonality of its characters gives equality of the unordered pairs, including repetitions. Conversely conjugation preserves every real observation.
\end{proof}


\section{Examples, scope and further questions}
The archive's two Python programs illustrate the algebra and use the standard library alone. In \texttt{verify.py}, sampling at stride six collapses two distinct nodes at a Fourier bin. Multiplicities two and one produce a divided coefficient of modulus $\sqrt7$, while two isolated bins give their actual multiplicities two and three. Another calculation recovers the four weights $(1/2,-1/4,1/4,1/2)$ from four prescribed isolated bins, with the degree bound and multiplicities supplied for this finite example. A calculation with the first two positive zeta ordinates checks the geometric-sum error for finite Fourier means. The off-circle examples in \texttt{check\_shells.py} are synthetic tests of the general proof.

Primitivity fixes the arithmetic object being identified. Inducing a character to a larger modulus introduces Euler factors whose zeros lie on $\Re s=0$, leaving the open-strip multiset unchanged. The complex phases in \eqref{eq:li} retain the orientation of a character. Real parts of the unfiltered coefficients agree for $\chi$ and $\bar\chi$. The normalization $A(1)=1$ supplies the known unit in the integer-rounding step. Allowing polynomial degrees to grow without restriction on each window would permit exact interpolation of every finite list of discrepancies.

Two questions arise from the proof. First, for bounded conductors, strides and filter degrees, what effective separation time follows from explicit zero counts, a lower bound on the distance between candidate filters, and controlled measurement error? Second, for automorphic $L$-functions of unknown degree, the leading zero-count term measures degree times stride. What extra spectral information separates those two integers?

\begin{thebibliography}{9}\small
\bibitem{repo} \texttt{lsernau/grh-lean}, \emph{GRH for Dirichlet $L$-functions, in Lean 4}, based on L. Alp\"oge's proof with Claude. \href{https://github.com/lsernau/grh-lean/tree/72fead102f663d5c7a407964a13dc0cf9d977207}{Repository revision \texttt{72fead102f66}}, 25 November 2026.
\bibitem{lean} Lean FRO and Mathlib Initiative, \href{https://leanprover-community.github.io/blog/posts/riemann-hypothesis/}{\emph{Lean verification of the Riemann Hypothesis}}, November 2026.
\bibitem{wiki} Wikipedia contributors, \href{https://en.wikipedia.org/wiki/Riemann_hypothesis}{\emph{Riemann hypothesis}}, section on analytic criteria equivalent to RH, accessed 26 November 2026.
\bibitem{li} X.-J. Li, \href{https://doi.org/10.1006/jnth.1997.2137}{The positivity of a sequence of numbers and the Riemann hypothesis}, J. Number Theory \textbf{65} (1997), 325--333.
\bibitem{lagarias} J. C. Lagarias, \href{https://doi.org/10.5802/aif.2311}{Li coefficients for automorphic $L$-functions}, Ann. Inst. Fourier \textbf{57} (2007), 1689--1740.
\bibitem{bennett} M. A. Bennett, G. Martin, K. O'Bryant and A. Rechnitzer, \href{https://arxiv.org/abs/2005.02989}{Counting zeros of Dirichlet $L$-functions}, Math. Comp. \textbf{90} (2021), 1455--1482.
\bibitem{banks} W. D. Banks, \href{https://arxiv.org/abs/2302.07073}{Comparing zeros of distinct Dirichlet $L$-functions}, arXiv:2302.07073, revised 2024.
\end{thebibliography}
\end{document}
